\section{Discussion and Limitations}\label{sec:discussion}

The results answer the three research questions with a qualified conclusion. For RQ1, \sysname{} is competitive with the label-rich reference on GraSec-IoT and leads on Unicorn Wget, but the supervised GNN remains strongest on StreamSpot and GraSec-IoT. For RQ2, topology-aware scoring coincides with large F1 differences relative to GraphCL on GraSec-IoT and Unicorn Wget, but only a $0.003$ difference on StreamSpot. For RQ3, \sysname{} balances precision and recall well on GraSec-IoT and Unicorn Wget, whereas GraphCL attains the highest StreamSpot recall. These outcomes argue for dataset-dependent benefits, not universal superiority.

The supervised GNN's advantage on StreamSpot and GraSec-IoT is fair and expected: it is optimized directly with both benign and malicious labels, while \sysname{} and GraphCL do not use attack labels during representation learning. When labeled training attacks resemble test attacks, supervision provides direct information for locating a decision boundary. Its lower and more variable recall on Unicorn Wget shows that this advantage did not transfer uniformly in our evaluation. The available results do not establish why; possible contributors include the smaller sample, differences in graph construction, and differences between labeled training and test behavior. We therefore avoid attributing the result to overfitting without further analysis.

Dataset construction also provides a cautious explanation for the differing topology-aware gains. StreamSpot contains provenance graphs of heterogeneous system activity, GraSec-IoT represents IoT communication snapshots, and Unicorn Wget captures a specific system-activity scenario. These graph types can encode attacks at different structural scales. Degree- and closeness-based summaries may add little when pooled message-passing features already separate the relevant StreamSpot behavior, but may add complementary information when anomalous communication or provenance changes connectivity more broadly. Confirming this interpretation would require attack-level structural analysis or controlled ablations beyond the present study.

From an ARTMAN perspective, \sysname{} addresses one aspect of resilient and trustworthy ML-driven security: maintaining an anomaly-detection capability when comprehensive attack labels are unavailable. Learning a benign reference and scoring deviations can reduce reliance on a fixed catalog of attacks and can accommodate changing deployment behavior through an updated benign memory. This is not a claim of adversarial robustness. The evaluation measures ordinary detection performance under the reported splits; it does not test evasion, poisoning, adaptive attackers, or certified robustness.

\subsection{Limitations and Threats to Validity}

Several limitations constrain the conclusions. First, only three datasets and their existing graph constructions are evaluated, so the results may not generalize to other networks, organizations, attack families, or graph-generation procedures. Second, the small StreamSpot difference is within the scale of run-to-run variation and should not be treated as a substantive improvement. Third, labeled malicious validation graphs are used to select the F1-maximizing threshold. Although attack labels do not train the representation, this choice is not fully label-free and may be difficult to reproduce when deployment-time attacks are unavailable.

Fourth, balanced evaluation sets and F1-selected thresholds do not reproduce the low attack prevalence of many operational environments. Precision and alert volume may change under realistic class imbalance. Fifth, the supervised and benign-only methods have different access to labels, so the supervised GNN is a label-rich reference rather than a strictly matched baseline. Sixth, graph truncation, augmentations, hyperparameters, and the selected degree and closeness filtrations may affect performance; the present experiments do not establish which design component causes each dataset-level difference. Finally, \sysname{} detects anomalies at graph level and does not localize the nodes, edges, or subgraphs responsible for an alert, limiting analyst-facing explanation.

Future work should evaluate additional attack families and deployment shifts, select thresholds without labeled attacks, test under realistic class imbalance, and study sensitivity to graph construction and topology choices. Attack localization and structural explanations would make alerts more actionable. Robustness to adversarial manipulation, poisoning, and adaptive evasion also requires separate experiments before \sysname{} can be characterized as adversarially robust.
